\chapter{Typed λ-terms and unification problems}\label{ch:lambda}

\Cref{ch:horn,ch:hh} did logic programming over \emph{first-order} terms. \Cref{ch:hh} added
universal quantifiers and implications to goals (\fohh), and with them eigenvariables: constants made
up during the search that no logic variable created earlier may capture. Chapter~4 of the book
\cite{MillerNadathur2012} keeps the logic and changes the \emph{terms}. First-order terms give way to
\emph{simply typed λ-terms}. This chapter studies those terms on their own and leaves the higher-order
logic built on them to \cref{ch:hoprog} \book{96}.

The change has three consequences, and each of them touches a different part of a miniKanren:
\begin{enumerate}[nosep]
  \item \emph{What equality means.} Two terms are equal when they are λ-convertible (α, β, η), not
        only when they are identical. This affects the term representation and the unifier.
  \item \emph{What a logic variable may be bound to.} A unification problem becomes a
        \emph{quantified formula}. Its prefix mixes ∀ and ∃, and the order of the quantifiers decides
        which symbols a variable may use. This affects the state, which has to track scope.
  \item \emph{How hard unification is.} There may be several, infinitely many, or no most general
        unifiers, and the problem is undecidable. This affects the search: unification has to become a
        stream of answers.
\end{enumerate}
The chapter follows the book's order. \Cref{sec:lambda-syntax,sec:lambda-conversion,sec:lambda-normal}
cover the calculus. \Cref{sec:lambda-unifprob,sec:lambda-raising,sec:lambda-solutions,sec:lambda-examples}
recast unification as quantified equalities. \Cref{sec:lambda-vocab} fixes the vocabulary that
\cref{ch:hounif} uses for the unification procedure, and \cref{sec:lambda-undecidable} proves
undecidability.

% ======================================================================================
\section{Syntax and types of λ-terms}\label{sec:lambda-syntax}

Typed λ-terms are built from typed constants and variables with two operations. The first is
\emph{application}, written as juxtaposition and associating to the left, exactly as for first-order terms.
The second is \emph{abstraction}, which is new: from an expression $t$ in which a variable $x{:}\alpha$
may occur free, it forms the function $\lambda(x{:}\alpha)\,t$ \book{97}. In λProlog this is written
\lp{(x:A)\ T}. The type annotation can be dropped (\lp{x\ T}), and the type is then inferred.

\paragraph{Parsing.} The body of an abstraction extends as far to the right as possible. So
application binds tighter than abstraction, and abstraction associates to the right \book{97}:
\begin{center}
\begin{tabular}{@{}ll@{}}
\toprule
mathematical & λProlog concrete syntax\\
\midrule
$\lambda f\,\lambda x\,(f\,(f\,(f\,x)))$ & \lp{f\x\ f (f (f x))}\\
$\lambda x\,(f\,(g\,\lambda y\,(h\,x\,y))\,x)$ & \lp{x\ f (g y\ h x y) x}\\
$\lambda x\,\lambda y\,x$ & \lp{x\y\ x}\\
\bottomrule
\end{tabular}
\end{center}
A backslash cannot occur inside an identifier, so no spaces are needed around it.

\begin{definition}[Types and terms]\label{def:lambda-syntax}
Types are generated from primitive types (sorts such as $\typ{i}$, and the type $\typ{o}$ of
propositions), type variables $A$, type constructors such as $\typ{list}\,A$, and the arrow
$\tau_1\to\tau_2$, which associates to the right. Terms are
\[
  t \;::=\; c \;\mid\; x \;\mid\; (t\;t) \;\mid\; \lambda(x{:}\tau)\,t .
\]
(The book introduces the grammar in prose \book{97--99}. This compact form is ours.)
\end{definition}

\paragraph{Typing.} The judgement $\Sig;\Gamma\Vdash t:\tau$ (in the book, with a double turnstile) says
that $t$ has type $\tau$ when $\Sig$ gives the types of the constants and $\Gamma$ those of the
variables. It is defined by the rules of Figure~4.1 \book{98}:
\begin{mathpar}
\inferrule{c:\sigma\in\Sig \\ \tau\triangleleft\sigma}{\Sig;\Gamma\Vdash c:\tau}
\and
\inferrule{x:\tau\in\Gamma}{\Sig;\Gamma\Vdash x:\tau}
\and
\inferrule{\Sig;\Gamma\Vdash g:\tau_1\to\tau_2 \\ \Sig;\Gamma\Vdash t:\tau_1}{\Sig;\Gamma\Vdash (g\;t):\tau_2}
\and
\inferrule{\Sig;\Gamma,x:\tau\Vdash t:\sigma}{\Sig;\Gamma\Vdash \lambda(x{:}\tau)\,t:\tau\to\sigma}\;(\dagger)
\and
\inferrule{\Sig;\Gamma\Vdash B:\tau}{\Sig;\Gamma\Vdash C:\tau}\;(\ddagger)
\end{mathpar}
Proviso $(\dagger)$: $x$ is not declared as a type or kind in $\Sig$ or $\Gamma$. Proviso $(\ddagger)$:
$B$ and $C$ differ only in the names of bound variables. $\tau\triangleleft\sigma$ means that $\tau$ is a
substitution instance of $\sigma$, so a polymorphic constant is instantiated afresh at every use
\book{97}. Compared with the first-order rules of \cref{ch:horn}, the restrictions on the types in
$\Sig$ and $\Gamma$ are gone: a constant may have any type, including a functional type that ends in
$\typ{o}$. There is one new rule, for abstraction. And there are \emph{no special rules for the
quantifiers}, because the quantifiers have become ordinary constants (see below) \book{97--98}.

\begin{definition}[Well-formed, closed, simply typed {\book{98}}]\label{def:lambda-wf}
A term $t$ is a \emph{well-formed λ-term over $\Sig$} if $\Sig;\Gamma\Vdash t:\tau$ is derivable for
some $\Gamma$ and $\tau$. It is \emph{closed} if this holds with $\Gamma$ empty, and \emph{open}
otherwise. It is \emph{simply typed} if neither its type nor the types of its subterms contain type
variables.
\end{definition}

When every bound variable carries its annotation, the type of a well-formed term is unique up to
renaming of type variables. When annotations are omitted, type inference gives each bound variable a
fresh type variable and refines it \enquote{in as minimal a way as possible}. The type found this
way is again unique up to renaming, the principal type \book{98--99}. The unification theory of
\cref{ch:hounif} is stated for simply typed terms.

\begin{example}[A typing derivation]\label{ex:lambda-typing}
The Church numeral \lp{f\x\ f (f x)} has type $(\typ{i}\to\typ{i})\to\typ{i}\to\typ{i}$. Write
$\Gamma=f{:}\typ{i}\to\typ{i},\,x{:}\typ{i}$:
\begin{mathpar}
\inferrule*[right=($\dagger$)]{
  \inferrule*[right=($\dagger$)]{
    \inferrule*[right=app]{
      \inferrule*{f:\typ{i}\to\typ{i}\in\Gamma}{\Sig;\Gamma\Vdash f:\typ{i}\to\typ{i}} \\
      \inferrule*[right=app]{\Sig;\Gamma\Vdash f:\typ{i}\to\typ{i} \\ \Sig;\Gamma\Vdash x:\typ{i}}{\Sig;\Gamma\Vdash f\,x:\typ{i}}
    }{\Sig;\Gamma\Vdash f\,(f\,x):\typ{i}}
  }{\Sig;f{:}\typ{i}\to\typ{i}\Vdash \lambda x\,f\,(f\,x):\typ{i}\to\typ{i}}
}{\Sig;\emptyset\Vdash \lambda f\,\lambda x\,f\,(f\,x):(\typ{i}\to\typ{i})\to\typ{i}\to\typ{i}}
\end{mathpar}
Read the derivation bottom up. Each abstraction rule $(\dagger)$ moves one binder into the context.
The application rule \textsc{app} then checks $f$ against a functional type twice, and the leaves are
instances of the variable rule (the two leaves above $f\,x$ are left implicit). Because the derivation
ends with an empty context, the term is closed. No constant was used, so this derivation works for
every $\Sig$, including one with no constant of type $\typ{i}$, a case that matters in
\cref{sec:lambda-church,sec:lambda-solutions}.
\end{example}

\subsection{Formulas are terms of type o}\label{sec:lambda-formulas}

A \emph{formula} is a well-formed term of type $\typ{o}$. The connectives are constants such as
$\wedge:\typ{o}\to\typ{o}\to\typ{o}$. The quantifiers become constants too
\book{99}:
\[
  \forall : (A\to\typ{o})\to\typ{o}, \qquad \exists : (A\to\typ{o})\to\typ{o}
  \qquad\text{(concretely \lp{pi} and \lp{sigma})},
\]
and $\forall x\,B$ abbreviates $\forall(\lambda x\,B)$, written $\forall_\tau x\,B$ when the type of $x$
matters. A quantified formula is typed with the ordinary rules. From $\Sig;\Gamma,x{:}\sigma\Vdash B:\typ{o}$
the abstraction rule gives $\lambda x\,B:\sigma\to\typ{o}$, the instance
$\forall:(\sigma\to\typ{o})\to\typ{o}$ of the constant gives the rest, and so $\forall(\lambda x\,B):\typ{o}$.
For example, applying \lp{pi} at type \lp{(list int -> o) -> o} to an abstraction of type
\lp{list int -> o} \book{99}:
\begin{lprolog}
y\ append (1::2::nil) y X.
pi y\ append (1::2::nil) y X.
\end{lprolog}
Since signatures are no longer restricted, logical symbols may now appear \emph{inside the arguments of
non-logical ones} \book{99--100}:
\begin{lprolog}
type foreach  (A -> o) -> list A -> o.
foreach (x\ x > 5, x < 9) (3::10::6::8::nil).
\end{lprolog}
Here the conjunction \lp{,} sits inside an argument of \lp{foreach}. Higher-order programming
(\cref{ch:hoprog}) rests on this possibility. One more fact matters a lot later. \emph{The logic has a
single binder, λ-abstraction}: ∀, ∃ and every object-level binder (\cref{ch:lambdatree}) are encoded
with it \book{100}.

\begin{kanren}[terms are data, goals are procedures]
In \micro{} a goal is a Scheme procedure from a state to a stream of states. A procedure is opaque to
the unifier: it cannot be walked, compared or reified. So a formula on which one wants to
\emph{unify}, for example to solve for a predicate variable, has to be represented as
\emph{data}, such as a record for $\wedge$ or for \lp{pi} applied to an abstraction, and then
interpreted by a meta-level relation. This is the design of \cref{ch:hoprog,ch:lkanren}. For
\lp{foreach} with a ground Scheme predicate, ordinary higher-order Scheme is enough. Note also that
the term language now needs a representation of λ itself, which Scheme closures cannot provide:
\cref{sec:lambda-normal,sec:lambda-unifprob} use de Bruijn terms \cite{deBruijn1972} (lists in the first, records in the second), and \cref{ch:design}
discusses the options.
\end{kanren}

% ======================================================================================
\section{λ-conversion: the equality of the logic}\label{sec:lambda-conversion}

An abstraction is meant to denote a function, and an application of an abstraction is meant to denote
the result of the function. The conversion rules make this intended meaning part of the equality.
They rely on substitution, which must not capture variables.

\begin{definition}[Free for; substitution {\book{100}}]\label{def:lambda-subst}
$t$ is \emph{free for} $x$ in $s$ if no free occurrence of $x$ in $s$ lies in the scope of an
abstraction that binds a free variable of $t$. If $t$ is free for $x$ in $s$ and both have the same
type, $\subst{s}{t}{x}$ is $s$ with every free occurrence of $x$ replaced by $t$.
\end{definition}
For example, $(f\,x)$ is free for $u$ in $\lambda w\,(g\,u\,w)$, but $(f\,w)$ is not: substituting it
would put its free $w$ under the binder $\lambda w$. Ignoring the proviso is logically unsound, for the
same reason as with quantified formulas in \cref{ch:hh}.

\begin{definition}[The rules of λ-conversion {\book{100--101}}]\label{def:lambda-conv}
The following rewrites may be applied to any subterm:
\begin{center}
\begin{tabular}{@{}lll@{}}
\toprule
rule & redex $\Rightarrow$ contractum & proviso\\
\midrule
α-rewriting & $\lambda x\,s\;\Rightarrow\;\lambda y\,\subst{s}{y}{x}$ & $y$ free for $x$ in $s$, $y$ not free in $s$\\
β-contraction & $(\lambda x\,s)\,t\;\Rightarrow\;\subst{s}{t}{x}$ & $t$ free for $x$ in $s$ (lifted below)\\
η-contraction & $\lambda x\,(s\,x)\;\Rightarrow\;s$ & $x$ not free in $s$\\
\bottomrule
\end{tabular}
\end{center}
The converses are β-expansion and η-expansion. \emph{α-conversion} is the reflexive, symmetric,
transitive closure of α-rewriting. \emph{β-reduction} is the reflexive, transitive closure of
α-conversion together with β-contraction, and \emph{β-conversion} is its symmetric, transitive closure.
η-reduction and η-conversion are defined in the same way. \emph{λ-conversion} is the transitive
closure of α-, β- and η-conversion.
\end{definition}

λ-conversion is an equivalence and a congruence, and \emph{it is the equality of higher-order logic}
\book{101}. Two consequences stand out. First, terms related by λ-conversion cannot be told apart in
the logic. Second, \emph{no specification can determine the name of a bound variable}, because that
name can always be changed by α-conversion without changing the term. \Cref{ch:lambdatree} builds
λ-tree syntax on this second fact.

\begin{example}[Four equal terms {\book{101}}]\label{ex:lambda-four}
\begin{lprolog}
x\y\ f (g x) y
X\Y\ f (g X) Y
x\ f (g x)
x\y\ f ((u\v\v) (2 + 3) (g x)) y
\end{lprolog}
The first two terms are α-convertible. The third is the first after an η-contraction:
\[\lambda y\,(f\,(g\,x)\,y)\;\Rightarrow\;f\,(g\,x).\]
A β-contraction of the fourth gives
\begin{lprolog}
x\y\ f ((v\v) (g x)) y
\end{lprolog}
and a second one gives back the first. Note that \lp{(2 + 3)} disappears entirely, because $u$ does
not occur in the body \lp{v\v}.
\end{example}

\begin{pitfall}[capital letters are not always logic variables]
In \lp{X\Y\ f (g X) Y} the tokens \lp{X} and \lp{Y} are \emph{bound} names, not logic variables. In
λProlog only \emph{free} capitalised tokens are logic variables. A Scheme reader of λProlog code has to
look at where a name is bound before deciding what it is.
\end{pitfall}

\paragraph{Liberalised substitution.} From now on $\subst{s}{t}{x}$ is defined even when $t$ is not free
for $x$ in $s$. First α-rename $s$ to some $s'$ in which $t$ is free for $x$, then substitute. The result
depends on the renaming only up to α, so modulo equality it is well defined. β-contraction then becomes
$(\lambda x\,s)\,t\Rightarrow\subst{s}{t}{x}$ with no proviso \book{101}. This is ordinary
capture-avoiding substitution.

\begin{tutorial}[capture-avoiding β, by hand]
Contract $(\lambda u\,\lambda w\,(g\,u\,w))\,(f\,w)$. The argument has a free $w$. A naive textual
replacement would give $\lambda w\,(g\,(f\,w)\,w)$, where the outer $w$ has been \emph{captured} and now
refers to the bound variable. This is a different function. Liberalised substitution first renames the
binder, $\lambda u\,\lambda w'\,(g\,u\,w')$, and then substitutes:
\[
  (\lambda u\,\lambda w\,(g\,u\,w))\,(f\,w)\;\Rightarrow\;\lambda w'\,(g\,(f\,w)\,w').
\]
With de Bruijn indices (\cref{sec:lambda-normal}) the renaming becomes index arithmetic: the argument
is shifted by one when it goes under a binder.
\end{tutorial}

% ======================================================================================
\section{Normal forms and the power of β}\label{sec:lambda-normal}

A \emph{β-redex} is a term $(\lambda x\,s)\,t$. An \emph{η-redex} is a term $\lambda x\,(t\,x)$ with $x$
not free in $t$. A term with no β-redexes is in \emph{β-normal form}. If it has no η-redexes either, it
is in \emph{λ-normal form} (also called βη-normal form) \book{101--102}.

\begin{theorem}[Normal forms {\book{102}}]\label{thm:lambda-nf}
Every term of the typed language has a λ-normal form, unique up to α-conversion. It is written
$\mathrm{λnorm}(s)$. It can be computed by contracting β-redexes until none is left and then
contracting η-redexes.
\end{theorem}

\begin{corollary}[Decidable equality {\book{102}}]\label{cor:lambda-eq}
Two terms of the same type are λ-convertible if and only if their λ-normal forms are α-equivalent.
\end{corollary}

\begin{remark}[Background the book does not spell out]
The book states \cref{thm:lambda-nf} as \enquote{well known}. In the literature it rests on two
classical properties of the simply typed λ-calculus. \emph{Strong normalisation} says that every
reduction sequence terminates, which gives existence. The \emph{Church–Rosser property} (confluence)
says that diverging reductions can be joined again, which gives uniqueness. The algorithm also
relies on a standard fact the book uses without proof: η-contracting a β-normal term creates no new
β-redex, so doing η after β is enough.
\end{remark}

\begin{pitfall}[equality is intensional]
Logical equivalences do not affect term equality. With $p,q:\typ{i}\to\typ{o}$, the term
$\lambda x\,(p\,x)\wedge(q\,x)$ is \emph{not} equal to $\lambda x\,(q\,x)\wedge(p\,x)$ \book{102}. Types
are inhabited by \emph{expressions}, not by sets. A syntactic unifier gets this right automatically. A
unifier that knew about commutativity of $\wedge$ would be wrong.
\end{pitfall}

\subsection{Computing with β: Church numerals}\label{sec:lambda-church}

The rule $(\lambda x\,s)\,t\Rightarrow\subst{s}{t}{x}$ can copy $t$ many times or throw it away. Even when
$s$ and $t$ are both normal, the result need not be, because $x$ may sit in head position
\book{102}. So normalisation is a model of computation. Let the sort $\typ{i}$ have \emph{no constants}.
Then the closed λ-normal forms of type $(\typ{i}\to\typ{i})\to\typ{i}\to\typ{i}$ are exactly
the \emph{Church numerals} $\ulcorner n\urcorner=\lambda f\,\lambda x\,(f^n\,x)$ \book{102}. Arithmetic
is then definable \book{103}:
\[
  \mathit{succ}=\lambda n\,\lambda f\,\lambda x\,f\,(n\,f\,x),\quad
  \mathit{add}=\lambda n\,\lambda m\,\lambda f\,\lambda x\,n\,f\,(m\,f\,x),\quad
  \mathit{mul}=\lambda n\,\lambda m\,\lambda f\,\lambda x\,n\,(m\,f)\,x .
\]
λProlog prints λ-normal forms, so it can run these functions (footnote: the ambient signature has the
sort \lp{i}) \book{103}:
\begin{lprolog}
?- N = ((n\m\f\x\ n (m f) x) ((f:i -> i)\x\ f (f x)) (f\x\ f (f x))).
N = f\x\ f (f (f (f x)))
?-
\end{lprolog}
The book adds that the functions definable this way are \enquote{limited to the polynomial ones}
\book{103}. (The phrasing is informal. The precise class, a result outside the book due to
Schwichtenberg and Statman, is the \emph{extended} polynomials, which also allow a test for zero.)

\paragraph{Blow-up.} Define the \emph{size} of a term as its number of applications. Normalising can
increase the size enormously. The book's series \book{104} is
\begin{lprolog}
(g\e\ e)           (e\f\ e (e f)) (f\x\ f (f x)).
(g\e\ g e)         (e\f\ e (e f)) (f\x\ f (f x)).
(g\e\ g (g e))     (e\f\ e (e f)) (f\x\ f (f x)).
(g\e\ g (g (g e))) (e\f\ e (e f)) (f\x\ f (f x)).
\end{lprolog}
with \lp{x : i}, \lp{f : i -> i}, \lp{e : (i -> i) -> i -> i} and
\lp{g : ((i -> i) -> i -> i) -> (i -> i) -> i -> i}. The \lp{g\e\} parts are Church numerals one type
level up. \lp{g\e\ g (g e)} has a \emph{fourth-order} type. The $(n{+}1)$-th term has size $n+6$, and its
normal forms are the numerals 2, 4, 16 and 256 \book{104}.

\begin{remark}[An erratum in the size formula]\label{rem:lambda-erratum-size}
The book says the $(n{+}1)$-th normal form has size $2^{2^{\cdot^{\cdot^{2}}}}$ with $n+1$ twos
\book{104}. The term \lp{e\f\ e (e f)} squares a numeral, and the outer numeral applies it $n$ times to
$\ulcorner 2\urcorner$. So the normal form is $\ulcorner 2^{2^n}\urcorner$, of size $2^{2^n}$: 2, 4, 16,
256, 65536. This matches the book's own \enquote{256} for the fourth term. A right-nested tower with
four twos is 65536, though, so the printed formula is right only if the tower is read
left-nested, $((2^2)^2)^{\cdots}$. The run below confirms the numbers. The point of the passage is
unaffected: sizes can explode (Statman \cite{Statman1979} showed that λ-conversion is not elementary
recursive, as the book's notes cite \book{116}).
\end{remark}

\begin{kanren}[de Bruijn terms and a β-normaliser in plain Scheme]
With de Bruijn indices \cite{deBruijn1972} a bound variable is the number of binders between its
occurrence and its binder. α-equivalent terms are then \emph{the same data structure}, and
capture-avoidance turns into the index arithmetic of \scm{shift}. The following illustration is plain
Scheme, not a proposed API (the design is in \cref{ch:design}). It converts named terms to de Bruijn
form and computes β-normal forms (file \texttt{norm-core.scm}):
\begin{scheme}
;; A tiny β-normaliser: de Bruijn inside, names outside.
;; Input: (λ (x y ...) body), applications (f a b ...), symbols.
(define (->db t env)                    ; named term -> de Bruijn
  (cond
    ((symbol? t) (let ((i (index-of t env))) (if i i t)))
    ((eq? (car t) 'λ)
     (let loop ((xs (cadr t)) (env env))
       (if (null? xs)
           (->db (caddr t) env)
           (list 'lam (loop (cdr xs) (cons (car xs) env))))))
    (else
     (let loop ((f (->db (car t) env)) (as (cdr t)))
       (if (null? as)
           f
           (loop (list 'app f (->db (car as) env)) (cdr as)))))))
(define (index-of x l)
  (let loop ((l l) (i 0))
    (cond ((null? l) #f)
          ((eq? (car l) x) i)
          (else (loop (cdr l) (+ i 1))))))
(define (shift d c t)                   ; add d to the indices >= c
  (cond ((integer? t) (if (>= t c) (+ t d) t))
        ((symbol? t) t)
        ((eq? (car t) 'lam) (list 'lam (shift d (+ c 1) (cadr t))))
        (else (list 'app (shift d c (cadr t)) (shift d c (caddr t))))))
(define (subst t j s)                   ; t[s/j], capture-free
  (cond ((integer? t) (if (= t j) s t))
        ((symbol? t) t)
        ((eq? (car t) 'lam)
         (list 'lam (subst (cadr t) (+ j 1) (shift 1 0 s))))
        (else (list 'app (subst (cadr t) j s) (subst (caddr t) j s)))))
(define (beta body arg) (shift -1 0 (subst body 0 (shift 1 0 arg))))
(define (nf t)                          ; β-normal form
  (cond ((or (integer? t) (symbol? t)) t)
        ((eq? (car t) 'lam) (list 'lam (nf (cadr t))))
        (else (let ((f (nf (cadr t))))
                (if (and (pair? f) (eq? (car f) 'lam))
                    (nf (beta (cadr f) (caddr t)))
                    (list 'app f (nf (caddr t))))))))
(define (size t)                        ; number of applications
  (cond ((not (pair? t)) 0)
        ((eq? (car t) 'lam) (size (cadr t)))
        (else (+ 1 (size (cadr t)) (size (caddr t))))))
(define (church t)                      ; λf λx f^n x -> n, else #f
  (and (pair? t) (eq? (car t) 'lam)
       (pair? (cadr t)) (eq? (car (cadr t)) 'lam)
       (let loop ((b (cadr (cadr t))) (n 0))
         (cond ((eqv? b 0) n)
               ((and (pair? b) (eq? (car b) 'app) (eqv? (cadr b) 1))
                (loop (caddr b) (+ n 1)))
               (else #f)))))
(define (run t) (church (nf (->db t '()))))
\end{scheme}
Running Church arithmetic and the blow-up series (file \texttt{church.scm}):
\begin{scheme}
(include "norm-core.scm")
(define two '(λ (f x) (f (f x))))
(define three '(λ (f x) (f (f (f x)))))
(define (show . xs) (display xs) (newline))
(show 'succ-3  (run `((λ (n f x) (f (n f x))) ,three)))
(show 'add-2-3 (run `((λ (n m f x) (n f (m f x))) ,two ,three)))
(show 'mul-2-2 (run `((λ (n m f x) (n (m f) x)) ,two ,two)))
(define (series n)        ; (g\e\ g^n e) (e\f\ e (e f)) (f\x\ f (f x))
  `((λ (g e) ,(let loop ((k n)) (if (= k 0) 'e `(g ,(loop (- k 1))))))
    (λ (e f) (e (e f)))
    ,two))
(do ((n 0 (+ n 1))) ((> n 4))
  (let* ((t (->db (series n) '())) (v (nf t)))
    (show 'n n 'size (size t) 'nf-size (size v) 'numeral (church v))))
\end{scheme}
\begin{answers}
(succ-3 4)
(add-2-3 5)
(mul-2-2 4)
(n 0 size 6 nf-size 2 numeral 2)
(n 1 size 7 nf-size 4 numeral 4)
(n 2 size 8 nf-size 16 numeral 16)
(n 3 size 9 nf-size 256 numeral 256)
(n 4 size 10 nf-size 65536 numeral 65536)
\end{answers}
The input sizes are $n+6$, as the book says, and the normal-form sizes are $2^{2^n}$, as
\cref{rem:lambda-erratum-size} says. For the unifier, the useful lesson is that a de Bruijn λ-term
built from Scheme records or lists is \emph{already} handled by \micro's structural \scm{=°} up to α.
What \scm{=°} cannot do is β: binding a logic variable that sits in head position creates a redex, so
some walk-and-normalise step has to be added (\cref{def:module-hnf}).
\end{kanren}

\paragraph{Why logic programs do not explode.} The book gives three reasons \book{104--105}. (1)
Integers and data structures are built-in or made from non-logical constructors, not Church-encoded.
(2) Bound variables usually have primitive (order-0) types. If every bound variable in a term has a
primitive type and $s,t$ are λ-normal, then $\subst{s}{t}{x}$ is λ-normal: it may contain many copies of
$t$, but there are no new redexes. (3) In the sublanguage $\Llam$ (\cref{ch:lambdatree}), β is
restricted to \emph{β\textsubscript{0}-contraction} $(\lambda x\,s)\,x\Rightarrow s$, which only renames,
and normalisation then \emph{shrinks} terms.

\begin{remark}[Order of a type]
The book uses \enquote{$n$-th order} without giving a formula. The usual definition, consistent with all
its uses, is $\mathrm{ord}(\tau)=0$ for primitive $\tau$ and
$\mathrm{ord}(\tau_1\to\tau_2)=\max(\mathrm{ord}(\tau_1)+1,\mathrm{ord}(\tau_2))$. Checks: the type
$(\typ{i}\to\typ{i})\to\typ{i}\to\typ{i}$ is called second-order \book{102}, and the type of
\lp{g\e\ g (g e)} fourth-order \book{104}. A unification problem whose variables have order at most
2 is called \emph{third-order} \book{116}.
\end{remark}

% ======================================================================================
\section{Unification problems as quantified equalities}\label{sec:lambda-unifprob}

Section~1.5 of the book (\cref{ch:intro}) defined a unification problem as a multiset $\{t_1=s_1,\dots,t_p=s_p\}$ and related it
to the formula $\forall y_1\dots\forall y_m\,\exists x_1\dots\exists x_n\,[t_1=s_1\wedge\dots\wedge
t_p=s_p]$, with the $y$'s the constants and the $x$'s the logic variables \book{105}. Following Miller
\cite{Miller1991,Miller1992}, the book now makes the \emph{formula} the official presentation and
allows any quantifier prefix.

\begin{definition}[Generalised unification problem {\book{105}}]\label{def:lambda-unifprob}
A unification problem is a formula
\[
  Q_1x_1\dots Q_nx_n\,[\,t_1=s_1\wedge\dots\wedge t_m=s_m\,]\qquad(n,m\ge 0,\ Q_i\in\{\forall,\exists\}),
\]
with $t_j$ and $s_j$ λ-terms of the same type. The bound variables $x_i$ may have types of any order.
\end{definition}

Why mix the quantifiers? The book gives two reasons \book{106}, each with a query that fails.

\emph{First, \fohh{} proof search creates mixed prefixes.} A \lp{pi} in a goal introduces an
eigenvariable \emph{after} the logic variables already present:
\begin{lprolog}
?- pi y\ X = y.
no
?-
\end{lprolog}
The problem is $\exists x\,\forall y\,[x=y]$. It fails because no instance of \lp{X} may contain the
eigenvariable that stands for $y$.

\emph{Second, binders in terms behave like inner universal quantifiers.}
\begin{lprolog}
?- (y\ X) = (y\ y).
no
?-
\end{lprolog}
The problem is $\exists x\,[(\lambda y\,x)=(\lambda y\,y)]$. The equation $(\lambda y\,t)=(\lambda y\,s)$
is logically equivalent to $\forall y\,[t=s]$, so this is again $\exists x\,\forall y\,[x=y]$, and it
fails for the same reason: \emph{no instance of \lp{X} may contain a variable bound by an abstraction
inside whose body \lp{X} occurs} \book{106}.

\begin{tutorial}[reading a quantifier prefix]
Every problem in this chapter can be read with one rule.
\begin{quote}
\emph{An existential variable may be instantiated only with a term built from symbols quantified to its
\textbf{left}} (together with fresh logic variables when one is content with a unifier rather than a
solution, \cref{sec:lambda-solutions}).
\end{quote}
Here is where each quantifier of a λProlog problem comes from.
\begin{itemize}[nosep]
\item \textbf{Constants of the signature} are outermost universals. They are to the left of everything,
      so every logic variable may use them.
\item \textbf{Free variables of the query} (\lp{X}) are existentials.
\item \textbf{A \lp{pi} in a goal} adds a universal \emph{at that point}, to the right of the logic
      variables that already exist. Variables created later, for example by backchaining inside the
      scope of the \lp{pi}, sit to its right and may use it (\cref{ch:hh}).
\item \textbf{A λ inside an equation} is an innermost universal around that equation.
\end{itemize}
Four readings follow, with $g:\typ{i}\to\typ{i}\to\typ{i}$:
\begin{center}
\begin{tabular}{@{}lll@{}}
\toprule
problem & what the ∃ may use & verdict\\
\midrule
$\forall y\,\exists x\,[x=y]$ & $y$ & solved by $x:=y$\\
$\exists x\,\forall y\,[x=y]$ & nothing that mentions $y$ & no solution\\
$\forall a\,\exists F\,[(F\,a)=(g\,a\,a)]$ & $g$ and $a$ & four solutions\\
$\exists F\,\forall a\,[(F\,a)=(g\,a\,a)]$ & $g$ only & one solution\\
\bottomrule
\end{tabular}
\end{center}
The last two rows are worked out in \cref{sec:lambda-examples}. The same mechanism explains the
second failing query. In \lp{(y\ X) = (y\ y)} the logic variable \lp{X} was created
\emph{outside} the binder, so it is to the left of the inner $\forall y$. It may be instantiated with
anything except a term that mentions that $y$.
\end{tutorial}

\paragraph{From ∀∃ to ∀∃∀.} In \fohc{} the natural class is ∀∃ problems: constants first, then logic
variables. With λ-terms the natural class is \textbf{∀∃∀} \book{106}. The outer ∀'s are the
constants. The inner ∀'s come from the binders of the equations: an equation at type $\tau\to\sigma$
is replaced by $\forall(w{:}\tau)$ around an equation at type $\sigma$, and those quantifiers are moved
out over the whole conjunction.

\begin{example}[Turning λ-equations into a ∀∃∀ prefix {\book{106--107}}]\label{ex:lambda-prenex}
Let $f:\typ{i}\to\typ{i}$ and $a:\typ{i}$, and consider
\[
  \exists G\,\exists H\,[(\lambda w\,f\,a)=(\lambda w\,f\,(G\,w))\;\wedge\;(\lambda w\,f\,w)=(\lambda w\,H\,(H\,w))],
\]
where $G,H:\typ{i}\to\typ{i}$. The transformation has three steps.
\begin{enumerate}[nosep]
\item Make the constants explicit:
  $\forall f\,\forall a\,\exists G\,\exists H\,[\dots]$.
\item Replace each λ-equation by a quantified one:
  \[\forall f\,\forall a\,\exists G\,\exists H\,[\forall w\,[(f\,a)=(f\,(G\,w))]\wedge\forall w\,[(f\,w)=(H\,(H\,w))]].\]
\item Prenex (one possibility):
  \[\forall f\,\forall a\,\exists G\,\exists H\,\forall w\,\forall w'\,[(f\,a)=(f\,(G\,w))\wedge(f\,w')=(H\,(H\,w'))].\]
\end{enumerate}
Moving quantifiers to the prefix can introduce dependencies that are not really part of the problem:
$\forall w$ and $\forall w'$ now both scope over the whole conjunction, although each was attached to one equation. The book calls these \enquote{logically harmless}
and assumes prenex form from now on \book{107}.
\end{example}
The example only illustrates the transformation. By a short check of ours (the book does not
discuss it), it has no solution. The first equation forces $G=\lambda w\,a$. But a closed $H$ over
$f,a$ is $\lambda z\,f^k z$ or $\lambda z\,f^k a$, so $H\,(H\,w')$ is $f^{2k}w'$ or $f^k a$, never $f\,w'$.

\begin{kanren}[miniKanren has only ∃, and binders leak]
\scm{fresh°} is an existential, and there is no universal. The second failing query shows what this
costs once binders are present. Encode \lp{(y\ X) = (y\ y)} with de Bruijn records (illustrative
records, file \texttt{scope.scm}):
\begin{scheme}
(import (aux kanren micro) (aux base) (aux stream))
;; de Bruijn records: α-equivalent terms are equal data
(define-record lam body)
(define-record dbi n)
(define (show x) (display x) (newline))
;; λx.λy.x and λu.λv.u are the same record tree
(show (°->list #f (fresh° (q) (=° (make-lam (make-lam (make-dbi 1)))
                                   (make-lam (make-lam (make-dbi 1))))
                               (=° q 'equal))))
;; ?- (y\ X) = (y\ y).   λProlog answers no; plain =° answers yes
(show (°->list #f (fresh° (q) (=° (make-lam q) (make-lam (make-dbi 0))))))
;; ∃x[x = x]: the answer is a unifier, not a closed witness
(show (°->list #f (fresh° (q) (=° q q))))
\end{scheme}
\begin{answers}
((λ () (quote equal)))
((λ () (make-record-instance (quote dbi) 0)))
((λ (α) α))
\end{answers}
The first line is what we want: α for free. The second is a \emph{wrong yes}. \scm{q} is bound to a
dangling index, which is exactly the forbidden capture of the inner $\forall y$. (The third line is
discussed in \cref{sec:lambda-solutions}.) For the prefix discipline, a miniKanren therefore needs
(i) eigenvariables, that is, fresh constants that are never bound, and (ii) a \emph{scope check} on
every binding: a variable may not receive a symbol quantified to its right. In \micro{} the monotone
\scm{vars-count} index already records the order in which things are created, which is almost the
prefix order. The open question is how a variable inherits the restriction of the variable it is
bound to. λProlog implementations answer it with \emph{universe levels}
\cite{NadathurJayaramanKwon1995,DunchevEtAl2015}, and \cref{ch:design} spells out the proposal.
\end{kanren}

% ======================================================================================
\section{Simplifying prefixes: raising}\label{sec:lambda-raising}

A prefix can alternate as often as it likes. Raising is the tool that brings every prefix to the
∀∃∀ shape, or even to ∃∀. It moves an existential to the left of a universal and pays for it by making
the existential's type higher-order.

\begin{proposition}[Raising {\book{107}}]\label{prop:lambda-raising}
Let $x:\tau$, $y:\sigma$ and $h:\sigma\to\tau$. Then
\[
  \forall y\,\exists x\,G\quad\text{is provable iff}\quad\exists h\,\forall y\,\subst{G}{(h\,y)}{x}.
\]
\end{proposition}
\begin{proof}[Proof (the book's argument)]
($\Rightarrow$) A proof of the left formula replaces $y$ by a new constant $c$ and $x$ by a term $t$,
which may contain $c$. Prove the right formula with $h:=\lambda c\,t$ and $y:=c$. Note that $c$ is not
free in $\lambda c\,t$, as the order $\exists h\,\forall y$ requires. ($\Leftarrow$) A proof of the right
formula replaces $h$ by some $s$ and $y$ by a new $c$, and $c$ cannot be free in $s$. Prove the left
formula with $y:=c$ and $x:=(s\,c)$.
\end{proof}

Applied repeatedly, raising moves every existential to the left of the universals before it. As a rule,
for a block of universals:
\[
  \forall y_1{:}\sigma_1\dots\forall y_k{:}\sigma_k\;\exists x{:}\tau\;G
  \;\rightsquigarrow\;
  \exists h{:}\sigma_1\to\dots\to\sigma_k\to\tau\;\forall y_1\dots\forall y_k\;\subst{G}{(h\,y_1\dots y_k)}{x}.
\]
The \enquote{cost} is that the type of the existential is \emph{raised} by the types of the universals it
crosses \book{108}.

\begin{tutorial}[raising, step by step]
\textbf{Step 0: a problem whose answer depends on order.} Let $g:\typ{i}\to\typ{i}\to\typ{i}$ and
$a:\typ{i}$ be constants, and compare
\[
  (\mathrm{P}_1)\;\forall w\,\exists z\,[z=g\,a\,w]
  \qquad\text{and}\qquad
  (\mathrm{P}_2)\;\exists z\,\forall w\,[z=g\,a\,w].
\]
By the reading rule, in P\textsubscript{1} $z$ may use $w$ and $z:=g\,a\,w$ solves it. In
P\textsubscript{2} $z$ may not use $w$, and no term without $w$ equals $g\,a\,w$, so P\textsubscript{2}
has no solution.

\textbf{Step 1: raise P\textsubscript{1}.} Move $\exists z$ across $\forall w$. The new variable
$h:\typ{i}\to\typ{i}$ replaces $z$ by $(h\,w)$:
\[
  \exists h\,\forall w\,[h\,w=g\,a\,w].
\]
Now $h$ comes first and may not mention $w$. But it is \emph{applied} to $w$, so it gets $w$ back as an
argument. The solution is $h:=\lambda c\,g\,a\,c$, and the check is $(\lambda c\,g\,a\,c)\,w\to_\beta g\,a\,w$.
Translating back with the proof of \cref{prop:lambda-raising} gives $z:=h\,w=g\,a\,w$, the solution of
Step~0. P\textsubscript{2} cannot be raised, because its existential is already on the left. The
difference between the two problems is now visible in the syntax: $h$ is applied to $w$ in one, and
$z$ stands alone in the other.

\textbf{Step 2: raise over constants as well (∃∀ form).} Constants are outer universals, so one can
also raise across them. Take $\forall a\,\exists F\,[(F\,a)=(g\,a\,a)]$ with $F:\typ{i}\to\typ{i}$.
Raising $F$ across $\forall a$ gives $F':\typ{i}\to\typ{i}\to\typ{i}$, with $F$ replaced by $(F'\,a)$:
\[
  \exists F'\,\forall a\,[(F'\,a\,a)=(g\,a\,a)].
\]
$F'$ may not use $a$. It must be $\lambda y\,\lambda x\,(g\,s_1\,s_2)$ with each $s_i\in\{x,y\}$, so there
are \emph{four} solutions. Mapped back by $F:=F'\,a$, they are the four solutions of the original
problem, which \cref{sec:lambda-examples} lists. For instance, $F'=\lambda y\,\lambda x\,g\,y\,x$ gives
$F=\lambda x\,g\,a\,x$.

\textbf{What to remember.} Raising is an exact logical equivalence. It replaces \enquote{may use the
universals to my left} by \enquote{is a closed function applied to those universals}. A logic variable
under binders, \lp{X} inside \lp{y\ ...}, can always be read as a raised variable applied to the bound
variables, \lp{(H y)}. When the arguments are \emph{distinct} bound variables, as in $h\,w$, the
application is a \emph{higher-order pattern} (\cref{ch:lambdatree,ch:hounif}), and pattern problems
have most general unifiers. In Step~2 the argument $a$ is repeated, $F'\,a\,a$, which is not a pattern,
and indeed there is no most general solution.
\end{tutorial}

\paragraph{Skolemization, the dual.} Skolemizing an \emph{assumption} $\forall x\,\exists y\,D$ with
$x:\sigma$ and $y:\tau$ gives $\forall x\,\subst{D}{(f\,x)}{y}$, where $f:\sigma\to\tau$ is a new
\emph{Skolem constant}. For a \emph{goal} the roles swap: $\exists x\,\forall y\,G$ becomes
$\exists x\,\subst{G}{(f\,x)}{y}$, that is, $\forall f\,\exists x\,\subst{G}{(f\,x)}{y}$ \book{108}. The
two techniques are dual in two senses. Skolemization moves an \emph{∃ to a smaller scope} and
introduces a new \emph{constant} of raised type. Raising moves a \emph{∀ to a smaller scope} and
introduces a new \emph{existential variable} of raised type. Logically they differ: \enquote{relating
unifiers for Skolemized unification problems to original, un-Skolemized unification problems is
problematic, especially when unification involves variables of higher-order type} \book{108}. Raising,
in contrast, is an equivalence (\cref{prop:lambda-raising}), which is why the book uses it to normalise
prefixes.

\begin{pitfall}[raising creates higher-order variables]
Raising a first-order problem gives a higher-order one: in Step~1, $z:\typ{i}$ became
$h:\typ{i}\to\typ{i}$. A first-order miniKanren with ∀-goals cannot implement scope by raising unless
it also has λ-terms and β. Prenexing (\cref{ex:lambda-prenex}) has a related cost: it may add spurious
dependencies. These are logically harmless, but a scope-tracking implementation becomes less
permissive than it needs to be.
\end{pitfall}

\begin{kanren}[raising versus levels]
Raising is how the book \emph{explains} logic variables under \lp{pi} and λ. Implementations usually
do not raise up front. Nadathur, Jayaraman and Kwon tag each variable and constant with a \emph{universe index}
\cite{NadathurJayaramanKwon1995}. A binding is allowed only if the term's constants have tags no larger
than the variable's, and the tags of the term's variables are lowered to match. ELPI uses the same
\emph{levels} \cite{DunchevEtAl2015}. Nadathur and Linnell \cite{NadathurLinnell2005} show how to erase
the prefix and raise \emph{on the fly}, only over the universals that actually occur. For \micro{},
the choice is between raising, which needs only terms (λ, application, β) but makes every existential
higher-order, and levels, which need a new per-variable piece of state and keep variables first-order.
The two must agree on the examples of this chapter. \Cref{ch:design} makes the choice.
\end{kanren}

% ======================================================================================
\section{Unifiers, solutions and empty types}\label{sec:lambda-solutions}

Once a problem is a formula, solving it should mean proving it. The connection is subtler than it
looks. The higher-order version of \rulename{$\exists$R} \book{108--109} is
\begin{mathpar}
\inferrule{\seq{\Sig}{\Prog}{\subst{B}{t}{x}} \\ \Sig;\emptyset\Vdash t:\tau}{\seq{\Sig}{\Prog}{\exists_\tau x\,B}}\;\rulename{$\exists$R}
\end{mathpar}
The witness $t$ must be \emph{closed}, because its context is empty. Proving $\exists x\,[x=x]$ with
$x:\typ{i}$ therefore needs a closed term of type $\typ{i}$, and whether one exists depends on
$\Sig$. In the same way, $\exists x\,\exists y\,\exists z\,[x=(f\,y)\wedge y=(f\,z)]$ with
$f:\typ{i}\to\typ{i}$ has the most general unifier $\{\langle X,f\,(f\,Z)\rangle,\langle Y,f\,Z\rangle\}$
in the sense of the book's Section~1.5 (\cref{ch:intro}). The formula, however, is provable only if some closed term of type
$\typ{i}$ exists \book{109}.

\begin{definition}[Legitimate substitution, unifier, solution {\book{109}}]\label{def:lambda-solution}
Consider the ∀∃∀ problem
$\forall u_1\dots\forall u_p\,\exists x_1\dots\exists x_n\,\forall w_1\dots\forall w_q\,[t_1=s_1\wedge\dots\wedge t_m=s_m]$
and let $\theta$ be a substitution for $x_1,\dots,x_n$.
\begin{itemize}[nosep]
\item $\theta$ is \emph{legitimate} if no $w_j$ occurs free in any $\theta(x_i)$.
\item A legitimate $\theta$ is a \emph{unifier} if $\theta(t_j)$ and $\theta(s_j)$ are λ-convertible for
      every $j$.
\item A unifier is a \emph{solution} if the only free variables in the terms $\theta(x_i)$ are among
      $u_1,\dots,u_p$.
\end{itemize}
A problem can have a unifier and no solution. \emph{It is a solution that makes the quantified formula
provable.}
\end{definition}

In the book's Section~1.5 (\cref{ch:intro}) \enquote{solution} and \enquote{unifier} meant the same thing. They have to be separated
now because, unlike first-order classical logic, higher-order intuitionistic logic does not assume that
every type is \emph{inhabited}. The Church numerals show why this matters. They are the closed terms of
type $(\typ{i}\to\typ{i})\to\typ{i}\to\typ{i}$ only if $\typ{i}$ has \emph{no} closed term $t$, since
otherwise $\lambda f\,\lambda x\,t$ would be a further, non-numeral inhabitant \book{109--110}.

\begin{example}[Unifier but no solution]
Let $\Sig$ contain only $u,v:\typ{i}\to\typ{i}$, so $\typ{i}$ is empty. The problem $\exists x\,[x=x]$
has the unifier $\{x\mapsto X\}$ with $X$ a fresh logic variable. It has no solution, because there is
no closed term of type $\typ{i}$. Add a constant $a:\typ{i}$ and $\{x\mapsto a\}$ becomes a solution.
\end{example}

\begin{kanren}[miniKanren answers are unifiers]
The third answer of \texttt{scope.scm} above, \scm{((λ (α) α))} for $\exists x\,[x=x]$, is a
\emph{unifier} in the sense of \cref{def:lambda-solution}. The reified \scm{α} is a fresh logic
variable, not a closed witness. This is also λProlog's practice: answers may contain variables, and
leftover flex–flex constraints (\cref{sec:lambda-vocab}). An untyped miniKanren cannot see empty types
at all. Obtaining \emph{solutions} would need an explicit generator of closed inhabitants of a type
over a signature, itself a proof-search relation. \Cref{sec:lambda-undecidable} shows that the gap
between unifiers and solutions is exactly where undecidability lives.
\end{kanren}

% ======================================================================================
\section{Example problems and their solutions}\label{sec:lambda-examples}

These examples \book{110--111} are the core intuition pump of the chapter. Each one shows a way in
which higher-order unification differs from the first-order kind.

\subsection{Quantifier order changes the answers}\label{sec:lambda-order}

Let $g:\typ{i}\to\typ{i}\to\typ{i}$ be a constructor of $\typ{i}$.
\begin{itemize}
\item $\forall a\,\exists F\,[(F\,a)=(g\,a\,a)]$ has \textbf{four} solutions:
      $F\in\{\lambda x\,g\,a\,a,\ \lambda x\,g\,x\,a,\ \lambda x\,g\,a\,x,\ \lambda x\,g\,x\,x\}$.
\item $\exists F\,\forall a\,[(F\,a)=(g\,a\,a)]$ has the \textbf{unique} solution $F=\lambda a\,g\,a\,a$.
\end{itemize}
\emph{Walk-through.} $F:\typ{i}\to\typ{i}$ is closed, so $F=\lambda x\,b$ with $b$ built from $g$, $x$ and
whatever $F$ may use. The β-normal form of $F\,a$ is $\subst{b}{a}{x}$, and it must be $g\,a\,a$. So $b$ is
$g\,s_1\,s_2$ where each $s_i$ either \emph{is} $a$ or \emph{becomes} $a$ under the substitution, that is,
$s_i\in\{a,x\}$. In the first problem $a$ is to the left of $F$, so both choices are allowed for each
position: $2\times 2=4$. In the second, $F$ may not mention $a$, only $s_i=x$ remains, and there is one
solution. (\Cref{sec:lambda-raising} reached the same count of four by raising.)

\begin{kanren}[a first-order simulation of both prefixes]
In this example $F$ is applied to a single argument and its body is first-order, so the problem can be
simulated in today's \micro. Represent $F$'s body as a first-order term in which \scm{x} stands for the
bound variable, and relate a body to its instance $\subst{b}{a}{x}$ with a relation \scm{inst°}. The
universal $a$ is the symbol \scm{a}. The ∃F∀a prefix, \enquote{$F$ may not mention $a$}, is then the
constraint \scm{(absent° 'a body)}. This is an illustration, not general higher-order unification (file
\texttt{order.scm}):
\begin{scheme}
(import (aux kanren micro) (aux base) (aux stream))
;; F = λx.body, body a first-order term over g, a and x = (dbi 0)
(define-record dbi n)
(set-record-printer! dbi (λ (r port) (display "x" port)))
(define x (make-dbi 0))
(define (show x) (display x) (newline))
;; (inst° body arg out): out = body[arg/x], the normal form of (F arg)
(define-relation (inst° body arg out)
  (cond°
    ((=° body x) (=° out arg))
    ((symbol° body) (=° out body))
    ((null° body) (null° out))
    ((fresh° (h t h* t*)
       (=° body (cons h t)) (=° out (cons h* t*))
       (inst° h arg h*) (inst° t arg t*)))))
;; ∀a ∃F [(F a) = (g a a)]: F's body may mention the constant a
(show (°->list/ground
        (fresh° (body) (inst° body 'a '(g a a)))))
;; ∃F ∀a [(F a) = (g a a)]: a comes after F, F's body must avoid it
(show (°->list/ground
        (fresh° (body) (absent° 'a body) (inst° body 'a '(g a a)))))
\end{scheme}
\begin{answers}
((g x x) (g a x) (g x a) (g a a))
((g x x))
\end{answers}
These are the book's four answers and its one answer. Two points carry over to the real design. First,
the scope check is a \emph{constraint} on what a variable may contain, like \scm{absent°}, and it is
attached at the moment the variable is created. Second, a relation that runs substitution
\emph{backwards} (\scm{inst°} with \scm{out} known) is what imitation and projection do in Huet's
procedure (\cref{ch:hounif}): each \scm{a} in the output is either imitated or projected from the
argument.
\end{kanren}

\subsection{Words over a two-letter alphabet}\label{sec:lambda-words}

Let the only constructors with target $\typ{i}$ be $u,v:\typ{i}\to\typ{i}$. Then $\typ{i}$ has no closed
terms, and the closed terms of type $\typ{i}\to\typ{i}$ are \emph{words}: $\lambda w\,w$ is the empty word
$\varepsilon$, $\lambda w\,u\,w$ is \enquote{u}, and $\lambda w\,u\,(v\,(u\,w))$ is \enquote{uvu}. Function
composition is concatenation \book{110}. Three problems, all with $F,G:\typ{i}\to\typ{i}$:
\begin{enumerate}
\item $\exists F\,\exists G\,[(\lambda w\,F\,(G\,w))=(\lambda w\,u\,(v\,(u\,w)))]$ has \textbf{four}
      solutions, the four ways of splitting uvu:
      $(F,G)\in\{(\mathrm{uvu},\varepsilon),(\mathrm{uv},\mathrm{u}),(\mathrm{u},\mathrm{vu}),(\varepsilon,\mathrm{uvu})\}$.
      Replacing the right-hand side by any other closed term asks for all ways of splitting that word.
\item $\exists F\,[(\lambda w\,u\,(F\,(u\,w)))=(\lambda w\,u\,(v\,(v\,(u\,w))))]$ has the \textbf{only}
      solution $F=\lambda w\,v\,(v\,w)$. The left-hand side is a \emph{pattern} that checks whether a word
      starts and ends with u.
\item $\exists F\,[(\lambda w\,u\,(F\,w))=(\lambda w\,F\,(u\,w))]$ is the word equation
      $\mathrm{u}\cdot F=F\cdot\mathrm{u}$. Its solutions are \textbf{infinitely many}:
      $F\in\{\lambda w\,w,\ \lambda w\,u\,w,\ \lambda w\,u\,(u\,w),\dots\}$, that is, $\mathrm{u}^k$ for every
      $k\ge 0$. They are closed and pairwise distinct, so none is an instance of another: \emph{there is
      no most general solution}.
\end{enumerate}

\begin{kanren}[words are lists, composition is append]
Because $\typ{i}$ is empty, a closed term of type $\typ{i}\to\typ{i}$ \emph{is} a word, and in this one case
higher-order unification becomes ordinary relational list programming (file \texttt{words.scm}):
\begin{scheme}
(import (aux kanren micro) (aux base) (aux stream))
;; closed terms of type i -> i over u, v : i -> i are words,
;; and composition is concatenation
(define-relation (append° r s rs)
  (cond°
    ((null° r) (=° s rs))
    ((fresh° (a d c) (cons° a d r) (cons° a c rs) (append° d s c)))))
(define (show x) (display x) (newline))
;; ∃F ∃G [(λw F (G w)) = (λw u (v (u w)))]
(show (μkanren-run (q -1 #t)
        (fresh° (F G) (append° F G '(u v u)) (=° q (list F G)))))
;; ∃F [(λw u (F (u w))) = (λw u (v (v (u w))))]
(show (μkanren-run (q -1 #t)
        (fresh° (uF) (append° '(u) q uF) (append° uF '(u) '(u v v u)))))
;; ∃F [(λw u (F w)) = (λw (F (u w)))]: five of infinitely many
(show (μkanren-run (q 5 #t)
        (fresh° (W) (append° '(u) q W) (append° q '(u) W))))
\end{scheme}
\begin{answers}
((() (u v u)) ((u) (v u)) ((u v) (u)) ((u v u) ()))
((v v))
(() (u) (u u) (u u u) (u u u u))
\end{answers}
The third query shows the right attitude toward problems without an mgu. A miniKanren does not need
a single answer: it returns a \emph{stream} of them, and \scm{take°} (here \scm{μkanren-run} with
$n=5$) cuts the stream. This is the shape a general higher-order unifier in \micro{} must have
(\cref{sec:lambda-undecidable}).
\end{kanren}

\subsection{Weak functionality}\label{sec:lambda-weak}

Let $\typ{i}$ have at least the constants $a,b,c,d:\typ{i}$. Then
$\exists F\,[(F\,a)=b\wedge(F\,c)=d]$ has \textbf{no solution}. The first equation is solved only by
$F=\lambda x\,b$: $F$ cannot produce $b$ by projecting its argument $a$, so it must be the constant
function. That $F$ fails the second equation \book{110--111}. A function that maps $a\mapsto b$ and
$c\mapsto d$ exists as a set-theoretic object, but it is \emph{not expressible} by these λ-terms, which
have no case analysis on constants. The book calls this \enquote{a relatively weak notion of
functionality}.

\begin{pitfall}[what a first-order intuition gets wrong]
\begin{itemize}[nosep]
\item \emph{Quantifier order is semantic}: $\forall a\,\exists F$ and $\exists F\,\forall a$ give four
      answers and one answer.
\item \emph{There may be several incomparable solutions} (four splittings of uvu), \emph{or
      infinitely many with no mgu} (u·F = F·u). A unifier returning one substitution, as
      \scm{μkanren-state-unify} does, cannot express this.
\item \emph{λ-terms are not arbitrary functions}: no case analysis, so $F\,a=b\wedge F\,c=d$ is
      unsolvable.
\end{itemize}
\end{pitfall}

% ======================================================================================
\section{Vocabulary for a unification procedure}\label{sec:lambda-vocab}

\Cref{ch:hounif} presents Huet's procedure \cite{Huet1975}. The book prepares it with some vocabulary
and a few structural observations \book{111--113}.

\begin{definition}[Binder, head, arguments, body {\book{111}}]\label{def:lambda-head}
Every β-normal term has the shape
\[
  \lambda x_1\dots\lambda x_n\,(h\;t_1\cdots t_p)\qquad(n,p\ge 0),
\]
where $h$ is a constant or a variable and every $t_i$ is β-normal. $x_1,\dots,x_n$ is the
\emph{binder}, $h$ the \emph{head}, $t_1,\dots,t_p$ the \emph{arguments}, and $(h\,t_1\cdots t_p)$ the
\emph{body}.
\end{definition}

\begin{definition}[βη-long normal form {\book{111}}]\label{def:lambda-long}
A term is in \emph{βη-long normal form} if it has the shape above, every $t_i$ is βη-long, and the body
has a non-functional type. For simply typed terms, that means a primitive type. Every β-normal form
can be η-expanded into an equal βη-long form.
\end{definition}
\begin{example}[η-expanding to long form {\book{111}}]
Let $f:(\typ{i}\to\typ{i})\to\typ{i}$. The term $f$ is β-normal but not long, since its body has
functional type. One η-expansion gives $\lambda w\,f\,w$ with $w:\typ{i}\to\typ{i}$. This is still not
long, because the argument $w$ has a non-primitive type. Expanding $w$ gives $\lambda w\,f\,(\lambda u\,w\,u)$
with $u:\typ{i}$, which is βη-long. Note that deciding when to stop requires knowing the \emph{types}.
\end{example}

\begin{definition}[Rigid and flexible {\book{111--113}}]\label{def:lambda-rigid}
In a problem over β-normal terms, an occurrence $\lambda\bar x\,(h\,t_1\cdots t_p)$ is \emph{rigid} if
$h$ is ∀-bound in the prefix (constants included) or belongs to the binder $\bar x$. It is
\emph{flexible} if $h$ is ∃-bound in the prefix. Equations are rigid–rigid, rigid–flexible,
flexible–rigid or flexible–flexible, and a problem is \emph{flexible–flexible} if all of its equations
are.
\end{definition}
The names are meant to suggest what substitution can do. Substitution never changes the head of a rigid
term: every instance has a β-normal form with the same head. It can change the head of a flexible term
\book{112}.

\paragraph{Observation 1: rigid heads must agree.} By α and η, two rigid terms of the same type can be
given the same binder: $\lambda\bar x\,(h\,t_1\cdots t_p)$ and $\lambda\bar x\,(k\,s_1\cdots s_q)$. Since
neither head can change, a unifier can exist only if $h$ and $k$ are the same symbol \book{112}. With $F,G$
logic variables and $c,d$ constants, none of the following has a unifier:
\begin{gather*}
  \{(\lambda x\,(d\,(c\,(F\,x))))=(\lambda x\,(c\,(d\,(G\,x))))\}\qquad
  \{(\lambda x\,(x\,(F\,x)))=(\lambda x\,(c\,(G\,x)))\}\\
  \{(\lambda x\,\lambda y\,(x\,(F\,x\,y)))=(\lambda x\,\lambda y\,(y\,(G\,x\,y)))\}
\end{gather*}
The first one, written precisely, is $\forall c\,\forall d\,\exists F\,\exists G\,[\lambda x\,(d\,(c\,(F\,x)))=\lambda x\,(c\,(d\,(G\,x)))]$.
Note that in the second and third problems one of the heads is a \emph{bound} variable, which counts as
rigid.

\paragraph{Observation 2: substitution cannot create captured occurrences.} If a binder variable has
no free occurrence in the body of a β-normal form, it has none in any instance either, because
substitution avoids capture. So
\[
  \{(\lambda x\,\lambda y\,x\,(F\,x\,y))=(\lambda x\,\lambda y\,G\,y\,y)\}
\]
has no unifier \book{112--113}. Every instance of the left side has head $x$. No instance of the right
side can contain the bound $x$: $G$ may be instantiated with a term containing a \emph{free} $x$, but
that is a different $x$.

\paragraph{Observation 3: flexible terms are limited only in their use of the binder.} A flexible term
can become almost any term of its type, \emph{except} in how it uses the binder variables. For
instance, $\lambda x\,\lambda y\,F\,(f\,x\,(G\,x))$ can become any $\lambda x\,\lambda y\,t$ with $x,y$ not
free in $t$: take $F:=\lambda z\,t$, which throws its argument away. Some instances also contain $x$,
such as $\lambda x\,\lambda y\,(f\,x\,x)$, obtained with $F:=\lambda z\,z$ and $G:=\lambda z\,z$. But no
instance contains $y$ \book{113}, because $y$ is passed to neither $F$ nor $G$, and by Observation~2
substitution cannot introduce it.

\paragraph{Observation 4: a flex–flex pair.} Unify $\lambda x\,\lambda y\,G\,x$ with
$\lambda x\,\lambda y\,H\,y$. Instances of the first ignore $y$, and instances of the second ignore $x$, so a
common instance is $\lambda x\,\lambda y\,t$ with neither variable free in $t$. The \enquote{most general}
unifier, derived in \cref{ch:hounif}, is $G:=\lambda w\,T$ and $H:=\lambda w\,T$, with $T$ a new variable
\book{113}. (Pattern unification calls this step \emph{pruning}, \cref{ch:hounif}.)

\begin{tutorial}[classifying heads in practice]
Take the prefix $\forall c\,\exists F\,\exists G$ and three occurrences:
\begin{center}
\begin{tabular}{@{}lll@{}}
\toprule
term & head & status\\
\midrule
$\lambda x\,(c\,(F\,x))$ & $c$, ∀-bound & rigid: every instance starts with $c$\\
$\lambda x\,(x\,(F\,x))$ & $x$, in the binder & rigid: every instance starts with the bound $x$\\
$\lambda x\,(F\,(c\,x))$ & $F$, ∃-bound & flexible: $F:=\lambda z\,z$ makes the head $c$\\
\bottomrule
\end{tabular}
\end{center}
A rigid–rigid equation either fails immediately or decomposes into equations between arguments
(\cref{ch:hounif}). A flex–rigid equation is where the search branches (imitate the rigid head, or
project an argument). A flex–flex equation always has a unifier, as in Observation~4, so Huet's
procedure \emph{postpones} it (\cref{ch:hounif}). It need not have a \emph{solution}, as
\cref{sec:lambda-undecidable} shows.
\end{tutorial}

\begin{kanren}[what the vocabulary means for the unifier]
\micro's \scm{=°} already performs \enquote{rigid–rigid}: two records or pairs with different
constructors fail at once, and the same constructor decomposes slot-wise. A de Bruijn representation
whose application node exposes head and arguments, as in \cref{def:lambda-head}, makes Observation~1
a constructor clash. Bound-variable heads (\scm{dbi} records) are rigid automatically. Everything about
\emph{flexible} heads is new: the unifier must recognise a logic variable in head position, normalise
after binding it, and either solve the equation deterministically (patterns, \cref{ch:lambdatree}) or
branch and delay it. Delayed flex–flex equations need a store in the state, next to the disequality
store, and \cref{ch:design} proposes one.
\end{kanren}

% ======================================================================================
\section{Undecidability}\label{sec:lambda-undecidable}

Derivability from \fohc{} programs is already undecidable, and so is unifiability of typed λ-terms
modulo λ-conversion \book{113}. The book proves the second fact by two reductions. It also stresses that
the specifications of later chapters produce \enquote{much tamer} problems.

\subsection{The Post correspondence problem}\label{sec:lambda-pcp}

\begin{definition}[PCP {\book{114}}]
Given $n\ge 1$ and two lists $s_1,\dots,s_n$ and $t_1,\dots,t_n$ of words over $\{u,v\}$, decide whether
there is a \emph{non-empty} sequence $i_1,\dots,i_k$ with $1\le i_j\le n$ such that
$s_{i_1}\cdots s_{i_k}=t_{i_1}\cdots t_{i_k}$. This problem is undecidable.
\end{definition}

Encode a word $r=r_1\cdots r_m$ as $\hat r=\lambda w\,(r_1\,(\cdots(r_m\,w)\cdots))$, as in
\cref{sec:lambda-words}. Concatenation is then $\lambda w\,(\hat r\,(\hat q\,w))$. Take a sort $\typ{i}$
\emph{with no constructors}, and introduce $u,v:\typ{i}\to\typ{i}$ by universal quantifiers
\book{114}:
\[
  \exists F\,\exists G\,\forall u\,\forall v\,
  [\,(F\,\hat s_1\cdots\hat s_n)=(F\,\hat t_1\cdots\hat t_n)\;\wedge\;(F\,u\cdots u)=\lambda w\,(u\,(G\,u\,w))\,],
\]
with $G:(\typ{i}\to\typ{i})\to\typ{i}\to\typ{i}$. For the lists $\mathrm{uv},\mathrm{u}$ and
$\mathrm{u},\mathrm{vu}$ \book{114}:
\[
  \exists F\,\exists G\,\forall u\,\forall v\,
  [\,(F\,(\lambda w\,u\,(v\,w))\,u)=(F\,u\,(\lambda w\,v\,(u\,w)))\;\wedge\;(F\,u\,u)=\lambda w\,(u\,(G\,u\,w))\,],
\]
where $F$ has the second-order type $(\typ{i}\to\typ{i})\to(\typ{i}\to\typ{i})\to\typ{i}\to\typ{i}$.

\begin{tutorial}[why the encoding is correct]
\textbf{Read the prefix.} $F$ and $G$ come \emph{before} $\forall u\,\forall v$, so they may not mention
$u$ or $v$. They can only rearrange their arguments. $\typ{i}$ has no constructors, so a closed $F$
must be $\lambda x_1\dots\lambda x_n\,\lambda w\,d$ with $d:\typ{i}$ built only from $x_1,\dots,x_n,w$.
The only possible shape is $d=x_{i_1}(\cdots(x_{i_k}\,w)\cdots)$ with $k\ge 0$ \book{114}.

\textbf{Read off the index sequence.} $F\,\hat s_1\cdots\hat s_n$ normalises to the word
$s_{i_1}\cdots s_{i_k}$, and the right side to $t_{i_1}\cdots t_{i_k}$. So the first equation says
exactly that $i_1,\dots,i_k$ solves the PCP, except in the degenerate case $k=0$, where both sides are
$\lambda w\,w$.

\textbf{Exclude $k=0$.} The second equation says that $F\,u\cdots u$ starts with $u$, which forces
$k\ge 1$. In the other direction, a PCP solution gives a unification solution directly
\book{114--115}. The argument extends to \emph{unifiers} and to types $\typ{i}$ with constructors
\book{115}.

\textbf{Check the instance.} For $\mathrm{uv},\mathrm{u}$ / $\mathrm{u},\mathrm{vu}$, the sequence
$1,2$ works: $\mathrm{uv}\cdot\mathrm{u}=\mathrm{u}\cdot\mathrm{vu}=\mathrm{uvu}$. This suggests
$F=\lambda x_1\,\lambda x_2\,\lambda w\,x_1\,(x_2\,w)$ and $G=\lambda y\,\lambda w\,y\,w$. $G$ may not mention
$u$ and has to obtain it from its argument. The book does not print this solution. The normaliser of
\cref{sec:lambda-normal}, with $u,v$ as free symbols, confirms both equations, and also the Diophantine
gadget of the next subsection (file \texttt{checks.scm}):
\begin{scheme}
(include "norm-core.scm")
(define (same? s t) (equal? (nf (->db s '())) (nf (->db t '()))))
(define (show . xs) (display xs) (newline))
;; PCP instance uv,u / u,vu; u and v are free symbols
(define F '(λ (x1 x2 w) (x1 (x2 w))))
(define G '(λ (y w) (y w)))
(show 'pcp (same? `(,F (λ (w) (u (v w))) u) `(,F u (λ (w) (v (u w)))))
           (same? `(,F u u) `(λ (w) (u (,G u w)))))
;; gadget for x = 1: λf λx N f (N f x) = λf λx N f (f x)
(define (numeral n)
  `(λ (f x) ,(let loop ((k n)) (if (= k 0) 'x `(f ,(loop (- k 1)))))))
(define (gadget N)
  (same? `(λ (f x) (,N f (,N f x))) `(λ (f x) (,N f (f x)))))
(show 'numerals-0..3 (map (lambda (n) (gadget (numeral n))) '(0 1 2 3)))
(show 'λf.λx.U (gadget '(λ (f x) U)))
\end{scheme}
\begin{answers}
(pcp #t #t)
(numerals-0..3 (#f #t #f #f))
(λf.λx.U #t)
\end{answers}
\end{tutorial}

Historically, Huet \cite{Huet1973} gave the PCP reduction, which shows that \emph{third-order}
unification is undecidable \book{116}.

\subsection{Hilbert's tenth problem}\label{sec:lambda-hilbert}

Hilbert's tenth problem asks for an algorithm that decides whether a system of Diophantine equations
(variables, numerals, $+$, $\times$ over $\mathbb N$) has a solution. No such algorithm exists
\book{115}. It is enough to treat equations of the forms $x=1$, $x+y=z$ and $x\times y=z$. Take a sort
$\typ{i}$ without constructors, so that the closed terms of $(\typ{i}\to\typ{i})\to\typ{i}\to\typ{i}$ are
exactly the Church numerals $\ulcorner n\urcorner$, and give every existential this type. Then
\book{115}:
\begin{enumerate}
\item The flex–flex problem $\exists N\,[(\lambda f\,\lambda x\,N\,f\,(N\,f\,x))=(\lambda f\,\lambda x\,(N\,f\,(f\,x)))]$
      has the unique solution $N=\ulcorner 1\urcorner$. With $N=\ulcorner n\urcorner$ the sides are
      $f^{2n}x$ and $f^{n+1}x$, so $2n=n+1$.
\item The flex–flex problem $\exists M\,\exists N\,\exists P\,[(\lambda f\,\lambda x\,N\,f\,(M\,f\,x))=(\lambda f\,\lambda x\,P\,f\,x)]$
      has the solution $\{(N,r),(M,s),(P,t)\}$ iff $r=\ulcorner n\urcorner$, $s=\ulcorner m\urcorner$ and
      $t=\ulcorner n+m\urcorner$ for some $n,m\ge 0$.
\item The flex–flex problem $\exists M\,\exists N\,\exists P\,[(\lambda f\,\lambda x\,N\,(M\,f)\,x)=(\lambda f\,\lambda x\,P\,f\,x)]$
      has the solution $\{(N,r),(M,s),(P,t)\}$ iff $r=\ulcorner n\urcorner$, $s=\ulcorner m\urcorner$ and
      $t=\ulcorner n\times m\urcorner$.
\end{enumerate}

\begin{remark}[An erratum in the addition gadget]
The book prints the left side of item 2 as $\lambda f\,\lambda x\,(N\,f)\,(M\,f)\,x$ \book{115}. This parses
as $((N\,f)\,(M\,f))\,x$, which is ill-typed, since $N\,f:\typ{i}\to\typ{i}$ is applied to
$M\,f:\typ{i}\to\typ{i}$. The intended term is the addition combinator of \cref{sec:lambda-church},
$\lambda f\,\lambda x\,N\,f\,(M\,f\,x)$, and that is the term shown above.
\end{remark}

\begin{theorem}[{\book{116}}]\label{thm:lambda-hilbert}
A system of such equations becomes a flex–flex problem with one existential variable per
$\mathbb N$-variable and one equation per equation. The problem has a \emph{solution} iff the system has
a solution in $\mathbb N$. The variables have order at most 2, so the problem is third-order, and the
construction extends to arbitrary Diophantine equations. Hence \textbf{it is undecidable whether a
third-order flex–flex problem has a solution.}
\end{theorem}

\paragraph{Unifiers versus solutions, once more.} The encoding \emph{always} has a unifier: pick a new
variable $U:\typ{i}$ and set every existential to $\lambda f\,\lambda x\,U$. Both sides of every gadget then
normalise to $\lambda f\,\lambda x\,U$. This term is not closed, so it is not a solution \book{116}. The run
above confirms both facts for gadget~1. Among $\ulcorner 0\urcorner,\dots,\ulcorner 3\urcorner$ only
$\ulcorner 1\urcorner$ works, and $\lambda f\,\lambda x\,U$ works too. So the undecidable question concerns
\emph{solutions}. Asking only for unifiers, a flex–flex problem is trivially unifiable, and this is why
Huet's procedure can leave flex–flex pairs for last. In the bibliographic notes \book{116--117} the book
credits this reduction to Miller \cite{Miller1992}. It differs from Goldfarb's \cite{Goldfarb1981}, which
already shows that \emph{second-order} unification, for unifiers, is undecidable.

\begin{kanren}[undecidability and the search]
Since unification is undecidable, no complete procedure terminates on every input. A complete
higher-order unifier is a \emph{semi-decision procedure}: a search, some of whose branches never
finish. In a miniKanren this is a familiar situation. Such a unifier would be a \emph{goal} that produces a
stream of unifiers. Huet's branching steps would become disjunctions, and \micro's \emph{interleaving}
\scm{orª} would explore them fairly, which a depth-first exploration does not. Flex–flex pairs, which always
have unifiers, would be left in a constraint store and reified with the answer, which then is a unifier and
not necessarily a solution. The proposal keeps only the deterministic part: pattern problems are solved by
one \scm{=°}, and non-patterns, flex–flex pairs included, wait in a store and are reported
(\cref{sec:module-wake}). Huet's branching is left to an optional search over that store, which needs
types (\cref{sec:roadmap-later-huet}).
\end{kanren}

\begin{remark}[Polymorphism]
All the formal results of the book's chapter concern \emph{simply typed} terms. Type variables are
allowed for convenience, for instance one declaration for \lp{pi}, but \emph{unification does not
instantiate type variables}: they behave like sorts, so the simply typed results still apply
\book{117}.
\end{remark}

% ======================================================================================
\section{Summary}\label{sec:lambda-summary}

\begin{itemize}
\item Terms are simply typed λ-terms (Figure~4.1). Formulas are terms of type $\typ{o}$, the quantifiers
      are constants $\forall,\exists:(A\to\typ{o})\to\typ{o}$ applied to abstractions, and λ is the
      \emph{only} binder.
\item The equality of the logic is λ-conversion (α, β, η). Bound names are unobservable, substitution
      avoids capture, and equality ignores logical equivalences.
\item Every term has a λ-normal form, unique up to α, so equality is decidable. β is powerful (Church
      numerals) and normal forms can be huge ($2^{2^n}$ in the book's series, see the erratum).
      Practical programs avoid this through built-in data, order-0 binders and $\Llam$.
\item A unification problem is a quantified formula $Q_1x_1\dots Q_nx_n\,[\bigwedge t_j=s_j]$.
      \emph{An existential may use only what is quantified to its left.} \lp{pi} and λ both introduce
      universals, which is why \lp{pi y\ X = y} and \lp{(y\ X) = (y\ y)} fail.
\item \emph{Raising} ($\forall y\,\exists x\,G\rightsquigarrow\exists h\,\forall y\,\subst{G}{(h\,y)}{x}$) is an
      exact equivalence that brings any prefix to ∀∃∀ or ∃∀, at the cost of higher-order types.
      Skolemization is its dual but does not preserve unifiers well.
\item A \emph{unifier} may leave free variables, while a \emph{solution} uses only the outer universals.
      Only solutions give proofs, and empty types make the difference real. miniKanren answers are
      unifiers.
\item Higher-order unification has several, infinitely many, or no solutions, and no mgu in general.
      Quantifier order changes the count (4 versus 1), and λ-terms cannot do case analysis.
\item Vocabulary for \cref{ch:hounif}: binder, head, arguments, body; βη-long form; rigid and
      flexible terms; rigid heads must agree; vacuous binders stay vacuous; flex–flex pairs always
      have unifiers.
\item Unifiability is undecidable (PCP, third order), and so is the existence of solutions of
      third-order flex–flex problems (Hilbert's tenth problem). Hence a miniKanren higher-order unifier
      must be a goal producing a fair stream of unifiers, with flex–flex pairs postponed.
\end{itemize}
